\SourceLeaf{344}

\Needspace{7\baselineskip}\begin{center}
[II.]

\textbf{[KREIS DURCH VIER PUNKTE.]}

[1.]
\end{center}

Der Kreis durch die Punkte, die den komplexen Größen $a$, $b$, $c$ mit den Konjugierten $A$, $B$, $C$ entsprechen, hat als Mittelpunkt den Punkt zu

\[
\begin{aligned}
&\frac{Aa(b-c)+Bb(c-a)+Cc(a-b)}{A(b-c)+B(c-a)+C(a-b)} \\
&= \frac{(Aa-Bb)(a-c)-(a-b)(Aa-Cc)}{(A-B)(a-c)-(a-b)(A-C)};
\end{aligned}
\]

einfacher, wenn $a=A=0$:

\[
\frac{(B-C)bc}{Bc-Cb}.
\]

\Needspace{7\baselineskip}\begin{center}
[2.]
\end{center}

Die Bedingungsgleichung dafür, dass vier Punkte $0$, $a$, $b$, $c$ mit den Konjugierten
$0$, $A$, $B$, $C$ auf ein und demselben Kreis liegen, lautet

\[
bcAB+caBC+abCA-abBC-bcCA-caAB = 0
\]
oder einfacher

\[
\frac{1}{aC}+\frac{1}{bA}+\frac{1}{cB}-\frac{1}{cA}-\frac{1}{aB}-\frac{1}{bC} = 0
\]
oder auch

\[
c(a-b)B(A-C) = b(a-c)C(A-B)
\]
etc.

\Needspace{7\baselineskip}\begin{center}
[3.]
\end{center}

Die Punkte, die zu dreien dieser Punkte in Bezug auf den vierten
reziprok sind, liegen auf ein und derselben Geraden.

Dieser Satz ist im allgemeineren Satz enthalten:

\begin{quote}
»Alle Punkte, die in Bezug auf einen gegebenen Punkt zu den
Punkten eines Kreises reziprok sind, liegen auf ein und demselben Kreis«.
\end{quote}
